\documentclass[10pt,a4paper]{amsart}
\usepackage[margin=19mm]{geometry}
\usepackage{amsmath,amssymb,mathtools}
\usepackage[scr=rsfs,cal=euler]{mathalfa}
\usepackage{xfrac}
\usepackage[unicode,hidelinks,bookmarks=false]{hyperref}
\newcommand{\ie}{i.e.\ }
\newcommand{\cf}{cf.\ }
\newcommand{\resp}{respectively\ }
\newcommand{\Subsection}{\subsection}
\providecommand{\nobreakdash}{\nobreak\hbox{-}}
\setlength{\emergencystretch}{2em}
\multlinegap0pt
\usepackage{bm}
\usepackage{bbm}
\usepackage{dsfont}
\usepackage{xspace}
\usepackage{cancel}
\usepackage{mathtools}
\usepackage{amsthm}
\makeatletter
\@ifundefined{theorem}{\newtheorem{theorem}{Theorem}}{}
\@ifundefined{lemma}{\newtheorem{lemma}[theorem]{Lemma}}{}
\makeatother
\title[On the structure of constacyclic codes over finite chain rin]{On the structure of constacyclic codes over finite chain rings}
\author{Vaishali Singh, Sucheta Dutt, Ridhima Thakral}
\date{Source: arXiv:2607.01771v1 (CC BY 4.0)}
\begin{document}
\maketitle

\noindent\emph{Source and licence.} On the structure of constacyclic codes over finite chain rings. Version arXiv:2607.01771v1 (CC BY 4.0), distributed under the Creative Commons Attribution 4.0 International licence, as recorded in the arXiv licence metadata for this exact version and shown on the abstract page https://arxiv.org/abs/2607.01771v1. Full rights notice: licenses/arxiv-2607.01771-CC-BY-4.0.txt.\par\smallskip

\noindent\textit{Editorial scope.} Counted unit: the Theorem whose source label is \texttt{thm:parttwo} in the article (source0.tex lines 216--225, source label \texttt{thm:parttwo}), delivered together with the article's own complete original proof of it (source0.tex lines 226--263, one \texttt{proof} environment of 38 lines). No mathematics is written, repaired or re-derived: statement and proof are the article's own text line for line. Required context retained uncounted: lem:first (lines 209--211). Every definition, display and label of those lines is retained unchanged. Omitted from this package and not redistributed: 1--208, 212--215, 264--461. Nothing inside a retained excerpt is omitted or altered: every retained body is delivered exactly as the article's lines are written, so no omission receipt of any kind accompanies this package. Citations: the retained excerpts carry no citation command at all, so no bibliography entry of the article is transcribed and no reference is redistributed. Reference adaptation: none was needed and none was made; every reference inside the delivered unit resolves inside it, so no unresolved reference or citation is produced. Printed slips kept literal: no printed word, symbol, hypothesis or number of the counted statement or of its proof was altered. Layout and class: the article's own class is \texttt{article} with packages including inputenc, fontenc, amsmath, amssymb, amsthm, mathtools, bm, xcolor, mathrsfs, hyperref, geometry, graphicx, booktabs, caption; the wrapper is the lane's pinned \texttt{amsart} editorial wrapper at 10pt on A4, which restates the theorem-like environments the retained text needs and adds the article's own macro definitions that the retained text needs (none), with extra wrapper packages (bm, bbm, dsfont, xspace, cancel, mathtools). The delivered numbering is the unit's own. Assets: no retained span contains a figure environment, a graphics-inclusion command, a PDF-page inclusion, a table or a TikZ picture; no image file is copied into this package, the package declares no asset dependency and no image file is redistributed with it. Licence: version arXiv:2607.01771v1 of this article is distributed under the Creative Commons Attribution 4.0 International licence, as recorded in the arXiv licence metadata for this exact version and shown on the abstract page https://arxiv.org/abs/2607.01771v1. A full rights notice is delivered at licenses/arxiv-2607.01771-CC-BY-4.0.txt. Characters: every retained excerpt is pure ASCII, so the delivered engine, XeLaTeX with its default Latin Modern text font, reports no missing character.
\begin{lemma}\label{lem:first}
    Let $\mathcal{C}=\langle f_0(x),f_1(x),\dots, f_s(x)\rangle$ be a $\lambda$-constacyclic code of length $\ell$ over $\mathcal{R}$, where $f_i(x),~0 \le i \le s$, be the polynomials as described above. Then, for every polynomial $q(x) \in \mathcal{C}$ such that $\deg(q(x))<n_i$, $q(x) \in \langle f_0(x),f_1(x), \dots, f_{i-1}(x) \rangle ~\text{for}~ 1 \le i \le s$.
\end{lemma}

\begin{theorem}
    Let $\mathcal{C}$ be a $\lambda$-constacyclic code generated by the set $\{ f_0(x),f_1(x), \dots , f_s(x) \}$, where $f_i(x),~ 0 \le i \le s$ are the polynomials as defined above. Then
    \begin{enumerate}
        \item $f_0(x)$, the minimum degree polynomial in $\mathcal{C}$, is uniquely determined.
        \item Each of the $f_i(x)$ is uniquely determined modulo $\langle f_0(x),f_1(x), \dots , f_{i-1}(x) \rangle$ for $1 \le i \le s$.
        \item $\gamma^{r_{i-1}-r_i}f_i(x) \in \langle f_0(x), f_1(x), \dots, f_{i-1}(x) \rangle$ for $1 \le i \le s$. \label{thm:parttwo}
        \item For $0 \le i \le s,~~ f_i(x)=\gamma^{r_i}h_i(x)$, where $h_i(x)$ is a monic polynomial.\label{thm:partthree}
        \item $h_0(x) \mid h_1(x) \pmod{\gamma^{\rho-r_0}}$, $h_{i-1}(x) \mid h_i(x) \pmod{\gamma^{r_{i-2}-r_{i-1}}}$ for $2 \le i \le s$, and $h_s(x) \mid (x^{\ell}-\lambda) \pmod{\gamma^{r_{s-1}-r_s}}$.
    \end{enumerate}
\end{theorem}

\begin{proof}
    \begin{enumerate}
        \item Suppose there exists another polynomial $g_0(x)$ such that $\deg(g_0(x))=n_0$ and the power of $\gamma$ in the leading coefficient of $g_0(x)$ is $r_0$. Then there exists a unit $u_0$ such that $f_0(x)-u_0g_0(x)$ is either $0$ or a polynomial in $\mathcal{C}$ of degree less than $n_0$. This is a contradiction. Therefore, $f_0(x)$ is uniquely determined.
        \item If there exists some $f_i(x),~1 \le i \le s$, that is not uniquely determined modulo $\langle f_0(x),f_1(x), \dots , f_{i-1}(x) \rangle$, then there exists $g_i(x)$ such that $\deg(g_i(x))=n_i$ and the power of $\gamma$ in the leading coefficient of $g_i(x)$ is $r_i$. Then there exists a unit $u_i$ such that $f_i(x)-u_ig_i(x)$ is either $0$ or a polynomial in $\mathcal{C}$ of degree less than $n_i$. By Lemma~\ref{lem:first}, $f_i(x)-u_ig_i(x) \in \langle f_0(x),f_1(x), \dots, f_{i-1}(x) \rangle$. It follows that $f_i(x)\equiv u_ig_i(x)$ modulo $\langle f_0(x),f_1(x), \dots, f_{i-1}(x) \rangle$. Hence, $f_i(x)$ is uniquely determined modulo $\langle f_0(x),f_1(x), \dots, f_{i-1}(x) \rangle$.
        \item Clearly, the polynomial $\gamma^{r_{i-1}-r_i}f_i(x)-x^{n_i-n_{i-1}}u_if_{i-1}(x) \in \mathcal{C}$, has degree less than $n_i$, where $u_i$ is a unit in $\mathcal{R}$. Therefore, by Lemma~\ref{lem:first}, $\gamma^{r_{i-1}-r_i}f_i(x)-x^{n_i-n_{i-1}}u_if_{i-1}(x) \in \langle f_0(x),f_1(x), \dots, f_{i-1}(x) \rangle$. It follows that $\gamma^{r_{i-1}-r_i}f_i(x) \in \langle f_0(x), f_1(x), \dots, f_{i-1}(x) \rangle$ for $1 \le i \le s$.
        \item The proof is carried out by induction on $i$. For $i=0$, let $f_0(x)= \gamma^{r_0}u_0x^{n_0}+ a_{n_0-1}x^{n_0-1}+ \dots+a_1x+a_0$, where $u_0$ is a unit in $\mathcal{R}$ and $a_i \in \mathcal{R},~ 0 \le i \le n_0-1$. If $a_k \not\equiv 0 \pmod{\gamma^{r_0}}$ for some $k, ~0 \le k \le n_0-1$, then $\gamma^{\rho-r_0}f_0(x)$ is a polynomial in $\mathcal{C}$ with degree less than $n_0$, which contradicts the fact that the degree of $f_0(x)$ in $\mathcal{C}$ is minimal. Therefore, $a_k \equiv 0 \pmod{\gamma^{r_0}}~ \forall~ 0 \le k \le n_0-1$ implying that $f_0(x)=\gamma^{r_0}h_0(x)$, where $h_0(x)$ is a monic polynomial. Thus, the result is true for $i=0$.

        Let us assume that the result is true for all $i \le k-1$ i.e. $f_i(x)=\gamma^{r_i}h_i(x)$, where $h_i(x)$ is a monic polynomial for all $0 \le i \le k-1$.

        Now we prove the result for $i=k$. We have $\gamma^{r_{i-1}-r_i}f_i(x) \in \langle f_0(x), f_1(x), \dots, f_{i-1}(x) \rangle$ by part~\ref{thm:parttwo}. Therefore, there exist polynomials $c_0(x),c_1(x),\dots,c_{k-1}(x)$ in $ \mathcal{R}[x]/ \langle x^{\ell}-\lambda \rangle$ such that
        \begin{align*}
            \gamma^{r_{k-1}-r_k}f_k(x) = \sum_{i=0}^{k-1}c_i(x)f_i(x)
            =\sum_{i=0}^{k-1}\gamma^{r_i}c_i(x)h_i(x)=\gamma^{r_{k-1}}\sum_{i=0}^{k-1}\gamma^{r_i-r_{k-1}}c_k(x)h_k(x)
        \end{align*}
        
        It follows that $\gamma^{\rho-r_k}f_k(x)=0$ and therefore, $f_k(x)=\gamma^{r_k}h_k(x)$, where $h_k(x)$ is a monic polynomial.
        
        Hence, $f_i(x)=\gamma^{r_i}h_i(x)$, where $h_i(x)$ is a monic polynomial for $0 \le i \le s$, by mathematical induction.
        
        \item From part~\ref{thm:parttwo}, we have $\gamma^{r_0-r_1}f_1(x) \in \langle f_0(x) \rangle$. Therefore, $\gamma^{r_0-r_1}f_1(x)=c_0(x)f_0(x)$, where $c_0(x) \in \mathcal{R}[x]/\langle x^{\ell}-\lambda \rangle$. This together with part~\ref{thm:partthree} implies that $\gamma^{r_0}(h_1(x)-c_0(x)h_0(x))=0$, which further implies that $h_1(x)-c_0(x)h_0(x)=0 \pmod{\gamma^{\rho-r_0}}$. Therefore, $h_0(x) \mid h_1(x) \pmod{\gamma^{\rho-r_0}}$.

        Again using part~\ref{thm:parttwo}, there exists $c_0(x),c_1(x), \dots, c_{i-1}(x) \in \mathcal{R}[x]/\langle x^{\ell}-\lambda\rangle$ such that
        \begin{align*}
            \gamma^{r_{i-1}-r_i}f_i(x)&=\sum_{k=0}^{i-1}c_k(x)f_k(x)
        \end{align*}
        Above equation together with part~\ref{thm:partthree} implies that
        \begin{align*}
            \gamma^{r_{i-1}}h_i(x)=\sum_{k=0}^{i-1}c_k(x)\gamma^{r_k}h_k(x)=\gamma^{r_{i-1}}c_{i-1}(x)h_{i-1}(x)+\sum_{k=0}^{i-2}c_k(x)\gamma^{r_k}h_k(x)
        \end{align*}
        It follows that
        \begin{align*}
            \gamma^{r_{i-1}}(h_i(x)-c_{i-1}(x)h_{i-1}(x))=\sum_{k=0}^{i-2}c_k(x)\gamma^{r_k}h_k(x)=\gamma^{r_{i-2}}\sum_{k=0}^{i-2}c_k(x)\gamma^{r_k-r_{i-2}}h_k(x)
        \end{align*}
           It follows that $\gamma^{\rho-(r_{i-2}-r_{i-1})}(h_i(x)-c_{i-1}(x)h_{i-1}(x))=0$, which implies that $h_i(x)-c_{i-1}(x)h_{i-1}(x)\equiv 0 \pmod{\gamma^{r_{i-2}-r_{i-1}}}$. Therefore, $h_{i-1}(x) \mid h_i(x) \pmod{\gamma^{r_{i-2}-r_{i-1}}}$ for $2 \le i \le s$.

        Now to prove that $h_s(x) \mid (x^{\ell}-\lambda) \pmod{\gamma^{r_{s-1}-r_s}}$, we see that $\gamma^{r_s}(x^{\ell}-\lambda) \in \mathcal{C}$ which implies that $\gamma^{r_s}(x^{\ell}-\lambda) \in \langle f_0(x),f_1(x), \dots, f_s(x) \rangle$. Now we will proceed in a similar manner as above to obtain the required result.
    \end{enumerate}
\end{proof}

\end{document}
